%% =====================================================================
%%  B4-T-22-03 -- what B4 contributes to "The Familiar Face"
%%  -------------------------------------------------------------------
%%  LAYER A: APPEND-ONLY. Written once at the entry's write, then left
%%  alone. If the source has more to say later, add a bullet -- do not
%%  rewrite. Material found ONLY here goes in the `unique` slot with
%%  @unique set to yes, which makes it undeletable (rule R10).
%% =====================================================================
%% @kind: source
%% @id: B4-T-22-03
%% @book: B4
%% @topic: T-22-03
%% @locator: ch. 5 (Liking), pp. 157--177 --- contact and cooperation, pp. 132--136
%% @status: distilled
%% @unique: yes
%% @integrated: yes
%% @updated: 2026-09-19

\dossier{B4}{T-22-03}{\bookshort{B4}}{ch. 5 (Liking), pp. 157--177 --- contact and cooperation, pp. 132--136}

\begin{distilled}
  \item For the most part, we like things that are familiar to us; the liking arrives without a reason attached.
  \item Subliminal exposure: faces flashed too briefly to recall still raised liking and persuasiveness in later interaction --- the exposure counter runs without memory.
  \item The mirror-image experiment: you prefer your reversed face, friends the true print --- each prefers the version seen most.
  \item Name familiarity wins elections: the Ohio candidate who adopted the state's political family name swept the race.
  \item Boundary condition: mere contact is not enough --- desegregation studies show contact under unequal or competitive conditions increasing prejudice; familiarity compounds liking only through cooperative contact.
\end{distilled}
\begin{unique}
  \item The subliminal-exposure finding, the mirror-image preference experiment, and the cooperative-contact boundary condition are unique to B4 in this corpus.
\end{unique}
